\documentclass[10pt, letterpaper]{article}

% ================= PACKAGES =================
\usepackage[
    ignoreheadfoot,
    top=2cm,
    bottom=2cm,
    left=2cm,
    right=2cm
]{geometry}

\usepackage{titlesec}
\usepackage{enumitem}
\usepackage{hyperref}
\usepackage{needspace}
\usepackage{iftex}

% ================= ATS / ENCODING (ALWAYS) =================
\usepackage[T1]{fontenc}
\usepackage[utf8]{inputenc}
\usepackage{charter}
\usepackage[final]{microtype}
\usepackage[french]{babel}

\input{glyphtounicode}
\pdfgentounicode=1

\pagestyle{empty}
\setcounter{secnumdepth}{0}
\setlength{\parindent}{0pt}
\renewcommand\labelitemi{-}

\titleformat{\section}{\bfseries\large}{}{0pt}{}[\titlerule]
\titlespacing{\section}{0pt}{6pt}{6pt}

\newenvironment{highlights}{
\begin{itemize}[leftmargin=12pt, itemsep=2pt]
}{
\end{itemize}
}

\hypersetup{
  colorlinks=true,
  urlcolor=black
}

% ================= DOCUMENT =================
\begin{document}

% ================= HEADER =================
{\fontsize{22}{22}\selectfont \textbf{Tamim Hmizi}} \\[4pt]
{\large Fondateur d'Axynoxia \;|\; Ingénieur Full Stack et Cloud DevOps} \\[6pt]
Ariana, Tunisie \;|\;
\href{mailto:tamimhmizi@icloud.com}{tamimhmizi@icloud.com} \;|\;
+216 54 886 068 \\
\href{https://linkedin.com/in/tamimhmizi}{linkedin.com/in/tamimhmizi} \;|\;
\href{https://github.com/tamim-hmizi}{github.com/tamim-hmizi} \;|\;
\href{https://tamim-hmizi.github.io/Portfolio/}{Portfolio}

\vspace{0.4cm}

% ================= PROFILE =================
\section{Profil Professionnel}
Ingénieur logiciel et \textbf{fondateur d'Axynoxia}, société d'ingénierie logicielle qui livre du développement sur mesure, du cloud et DevOps, de l'automatisation par IA, des intégrations et du e-commerce pour des clients en Europe et en Tunisie.
Je prends en charge la livraison de bout en bout : j'architecture le produit, je le développe full-stack, je le mets en production et j'exploite l'infrastructure qui le fait tourner.
\textbf{Plus de 6 ans d'expérience pratique} sur des plateformes enterprise, du SaaS multi-tenant, des applications mobiles et du e-commerce en production --- avec une expertise approfondie en Kubernetes, Terraform, automatisation CI/CD et architecture cloud-native sur Azure, AWS et OpenStack.
Réalisations : plateformes enterprise utilisées par \textbf{plus d'un million d'utilisateurs} dans des secteurs réglementés, infrastructures cloud hybrides souveraines pour un Partenaire Or Microsoft CSP, et automatisation apportant \textbf{70\% de réduction du temps de déploiement} et \textbf{40\% de réduction des tâches manuelles}.

% ================= SKILLS =================
\section{Compétences Techniques}
\textbf{Cloud et Plateformes :} Azure (Stack Hub, AKS, Public), AWS (EKS), GCP, OpenStack, DigitalOcean, Cloud Hybride, IaaS, PaaS \\
\textbf{Conteneurs et Orchestration :} Docker, Docker Compose, Kubernetes (AKS, EKS, K8s), Helm \\
\textbf{DevOps et Automatisation :} Terraform, Ansible, Jenkins, GitHub Actions, GitLab CI, Pipelines CI/CD, Infrastructure as Code (IaC), Nexus, Traefik, Portainer \\
\textbf{Monitoring et Observabilité :} Prometheus, Grafana, ELK Stack, CloudWatch, Exporteurs Python personnalisés, Alerting et Visualisation de métriques \\
\textbf{Sécurité et Qualité :} Pare-feux Fortinet, Segmentation réseau, IAM, SonarQube, Trivy, MFA, OAuth2, RBAC, Chiffrement des données \\
\textbf{Sauvegarde et Stockage :} Commvault, MinIO (Stockage objet), Redis (Cache), Stratégies de sauvegarde SQL et NoSQL \\
\textbf{Systèmes et Identité :} Linux (Ubuntu, CentOS), Windows Server, Active Directory, Microsoft 365, WSL \\
\textbf{Langages de Programmation :} TypeScript, JavaScript, Python, Java, C\#, PHP, Bash/Shell, SQL, C++, Assembleur \\
\textbf{Architectures Backend :} NestJS, FastAPI, Express.js, Django, Flask, Spring Boot, Node.js, APIs RESTful et WebSocket \\
\textbf{Architectures Frontend :} React (Vite, Next.js), Angular, Vue.js, Vuex, Tailwind CSS v4, Chakra UI, shadcn/ui \\
\textbf{Bases de données :} PostgreSQL, MongoDB, MySQL, Redis, Prisma, Mongoose, Sequelize, BullMQ (Files de traitement)

% ================= EXPERIENCE =================
\section{Expérience Professionnelle}

\textbf{Fondateur et Ingénieur Logiciel --- Axynoxia} \hfill Janv. 2026 -- Présent \\
Tunis, Tunisie
\begin{highlights}
\item Fondation et direction d'une \textbf{société d'ingénierie logicielle} au service de clients en Europe et en Tunisie --- développement sur mesure, cloud et DevOps, IA et automatisation intelligente, intégration de systèmes et e-commerce.
\item \textbf{Prise en charge du cycle de livraison complet} sur chaque mission : cadrage, architecture, développement full-stack, provisionnement cloud, CI/CD, observabilité et exploitation en production.
\item Mise en production et exploitation de systèmes réels : \textbf{plateforme de billetterie à forte concurrence}, \textbf{application mobile Flutter publiée sur l'App Store et Google Play}, \textbf{CMS immobilier} et \textbf{boutiques Shopify} pour des marques de distribution.
\item Développement de \textbf{produits SaaS multi-tenant} en interne --- ERP modulaire et plateforme de billetterie enterprise --- sur NestJS, Prisma, PostgreSQL, Redis et Traefik.
\item Gestion directe de la relation client, du cadrage technique et de la livraison commerciale, en \textbf{arabe, français et anglais}.
\end{highlights}

\textbf{Ingénieur Full Stack --- Bassetti Group (TEEXMA)} \hfill Avr. 2026 -- Présent \\
Télétravail
\begin{highlights}
\item Développement full stack sur des applications enterprise en \textbf{.NET}, \textbf{Angular}, \textbf{JavaScript}, \textbf{HTML} et \textbf{CSS} sur la plateforme \textbf{TEEXMA} --- une solution no-code collaborative de gestion des données techniques et scientifiques.
\item Maintenance et extension des \textbf{modules legacy Delphi} constituant le moteur de traitement central de TEEXMA.
\item Conteneurisation des charges applicatives avec \textbf{Docker} et contribution aux workflows d'automatisation et \textbf{DevOps}.
\item Contribution à un produit utilisé par \textbf{plus d'un million d'utilisateurs} dans des secteurs réglementés --- Aérospatial, Énergie, Automobile, Chimie et Pharmaceutique --- pour des clients tels qu'\textbf{Airbus}, \textbf{Safran}, \textbf{Framatome} et \textbf{Saint-Gobain}.
\end{highlights}

\textbf{Consultant Infrastructure --- RFC (Partenaire Or Microsoft CSP)} \hfill Oct. 2025 -- Mai 2026 \\
Ariana, Tunisie
\begin{highlights}
\item \textbf{Conception} et optimisation d'\textbf{infrastructures cloud hybrides en production} pour un acteur majeur du \textbf{Cloud Solution Provider (CSP) Microsoft}, intégrant \textbf{Azure Stack Hub} aux services Azure publics.
\item \textbf{Architecture} d'environnements cloud souverains pour des clients enterprise, garantissant la résidence des données et la conformité via des architectures certifiées Microsoft.
\item \textbf{Orchestration} de la sécurité et du réseau à l'échelle enterprise avec \textbf{Fortinet SD-WAN} et logique \textbf{Cisco} pour protéger une infrastructure multi-tenant.
\item \textbf{Administration} des opérations de sauvegarde critiques avec \textbf{Commvault}, mise en place de la validation automatisée des restaurations et rapports d'analyse des incidents.
\item \textbf{Automatisation} des tâches récurrentes et des correctifs de sécurité sur les parcs Linux et Windows Server, réduisant la charge manuelle de \textbf{40\%}.
\item \textbf{Animation} d'ateliers techniques de haut niveau en tant qu'\textbf{Expert Infrastructure (SME)}, accompagnant les clients sur la sécurité cloud-native et la gouvernance Microsoft 365.
\end{highlights}

\textbf{Ingénieur DevOps Plateforme --- Projet de Fin d'Études (PFE), RFC} \hfill Fév. 2025 -- Août 2025
\begin{highlights}
\item Conception et implémentation d'une \textbf{plateforme SaaS DevOps-as-a-Service} pour automatiser l'analyse, le déploiement et le monitoring des applications.
\item Développement d'une \textbf{application monolithique centrale} avec \textbf{React}, \textbf{Express.js}, \textbf{FastAPI} et \textbf{MongoDB}, avec séparation claire entre frontend, passerelle API et services backend.
\item Intégration d'un \textbf{moteur d'analyse basé sur l'IA (LLM)} pour inspecter les dépôts Git, la structure applicative, les dépendances et les exigences d'exécution.
\item Implémentation d'une \textbf{logique de déploiement multi-architecture}, permettant une prise de décision automatisée entre :
\begin{itemize}[leftmargin=16pt, itemsep=2pt]
  \item \textbf{Déploiements sur Machine Virtuelle} pour les applications monolithiques.
  \item \textbf{Déploiements Kubernetes (AKS)} pour les architectures microservices.
\end{itemize}
\item Développement de deux \textbf{applications e-commerce MERN} comme cas d'usage réels de déploiement :
\begin{itemize}[leftmargin=16pt, itemsep=2pt]
  \item Une \textbf{application e-commerce monolithique} déployée sur \textbf{Machine Virtuelle}, sélectionnée par le moteur IA selon la structure du dépôt.
  \item Un \textbf{backend e-commerce microservices} déployé sur \textbf{Kubernetes (AKS)}, recommandé par le moteur IA en raison de la décomposition des services.
\end{itemize}
\item Automatisation des \textbf{déploiements VM} avec \textbf{Terraform} pour le provisionnement, \textbf{Ansible} pour la configuration et \textbf{Docker / Docker Compose} pour la livraison applicative.
\item Implémentation des \textbf{déploiements microservices AKS} avec \textbf{AKS Engine}, \textbf{kubectl} et des \textbf{manifestes YAML Kubernetes} par microservice.
\item Conception et implémentation de \textbf{pipelines CI/CD Jenkins} couvrant les phases de build, test, qualité, sécurité, gestion des artefacts et déploiement.
\item Intégration de \textbf{SonarQube} pour l'analyse de la qualité du code et \textbf{Nexus} pour la gestion des artefacts.
\item Implémentation d'une \textbf{observabilité et d'un monitoring complets} avec \textbf{Prometheus} et \textbf{Grafana}.
\item Réduction du temps de provisionnement et de déploiement d'environ \textbf{70\%} grâce à l'automatisation.
\end{highlights}

\textbf{Ingénieur Cloud et Full Stack --- Stage, RFC} \hfill Juil. 2024 -- Août 2024
\begin{highlights}
\item Développement d'une \textbf{application web cloud-native} pour la \textbf{gestion des stagiaires}, couvrant l'intégration, le suivi et les workflows administratifs.
\item Déploiement de l'application sur \textbf{AWS EKS (Kubernetes)} dans une architecture entièrement conteneurisée.
\item Conception et implémentation de \textbf{pipelines CI/CD bout-en-bout} avec \textbf{Jenkins}, \textbf{Docker}, \textbf{Terraform} et \textbf{GitHub Actions}.
\item Implémentation des \textbf{contrôles de sécurité IAM} et du \textbf{monitoring applicatif et infrastructure} avec \textbf{Prometheus} et \textbf{Grafana}.
\end{highlights}

\textbf{Développeur Web --- Stage, ESPRIT} \hfill Juil. 2023 -- Août 2023
\begin{highlights}
\item Conception et développement d'une \textbf{plateforme web de scraping de sites éducatifs officiels} pour collecter et agréger des cours en ligne.
\item Implémentation de \textbf{modules de scraping Python} pour extraire les métadonnées des cours : titre, fournisseur, durée et niveau de difficulté.
\item Normalisation et stockage des données extraites dans une base de données structurée pour une recherche et un filtrage efficaces.
\item Développement de services backend pour classer et recommander les \textbf{meilleurs cours selon les critères de recherche}.
\item Création d'une interface web responsive permettant la recherche, la comparaison et la découverte de cours issus de sources fiables.
\item Garantie de la fiabilité des données en ciblant uniquement des \textbf{fournisseurs de cours officiels et vérifiés}.
\end{highlights}

\textbf{Stagiaire Ingénieur Réseau (C-NOC) --- Tunisie Telecom} \hfill Août 2021 -- Sept. 2021
\begin{highlights}
\item Travail au sein du Centre d'Opérations Réseau B2B national.
\item Surveillance des liens réseau enterprise et gestion des incidents de connectivité.
\item Assistance au dépannage des problèmes de routage et de disponibilité des services.
\end{highlights}

% ================= PROJECTS =================
\section{Projets}

\textbf{1. Plateforme de Monitoring Enterprise --- Azure Stack Hub (Production RFC)}
\begin{highlights}
\item \textbf{Vue d'ensemble :} Surveillance en temps réel de la santé de l'infrastructure Azure Stack Hub, du fabric de stockage et des performances des VM locataires (IOPS/Calcul).
\item \textbf{Conception} d'un écosystème de monitoring industriel pour des \textbf{architectures cloud hybrides} à l'aide d'exporteurs \textbf{Dockerisés} spécialisés.
\item \textbf{Développement} de micro-exporteurs haute performance en \textbf{Python/Flask} collectant la télémétrie depuis les \textbf{APIs Admin Azure Stack Hub} via l'authentification \textbf{OAuth2 par Principal de Service}.
\item \textbf{Architecture} de clusters de scraping \textbf{Prometheus} centralisés et de heatmaps multi-dimensionnelles \textbf{Grafana} offrant une visibilité sub-minute sur la santé des régions, du fabric de stockage et des IOPS disque des VM.
\item \textbf{Implémentation} de seuils d'alerting automatisés et de métriques d'observabilité avancées pour minimiser les temps d'arrêt en environnements cloud souverains.
\end{highlights}

\textbf{2. Axynoxia ERP --- Plateforme SaaS Multi-Tenant}
\begin{highlights}
\item \textbf{Vue d'ensemble :} Suite de gestion d'entreprise centralisée avec outils CRM, inventaire et comptabilité modulaires pour organisations multi-tenant.
\item \textbf{Architecture} d'un \textbf{monorepo SaaS ERP multi-tenant} haute performance avec \textbf{NestJS}, \textbf{Prisma}, \textbf{PostgreSQL} et \textbf{Redis}.
\item \textbf{Implémentation} de standards de sécurité industriels : \textbf{MFA (QR-based)}, \textbf{RBAC} et chiffrement \textbf{Argon2} pour les données métier sensibles.
\item \textbf{Orchestration} d'une infrastructure résiliente avec \textbf{Docker Compose} et \textbf{Traefik} pour la découverte automatique des services et le routage sécurisé des tenants.
\item \textbf{Développement} d'un système de modules auto-découverts permettant aux plugins métier (CRM, Inventaire) de s'intégrer au noyau sans modification du code.
\end{highlights}

\textbf{3. Tickety --- SaaS de Billetterie Haute Concurrence}
\begin{highlights}
\item \textbf{Vue d'ensemble :} Gestion d'événements à grande échelle, réservation de billets, transactions portefeuille et validation sur site pour les organisateurs.
\item \textbf{Architecture} d'un \textbf{monorepo SaaS distribué} (pnpm) avec \textbf{React}, \textbf{NestJS} et \textbf{Prisma}, supportant des milliers de transactions de réservation simultanées par seconde.
\item \textbf{Conception} de schémas \textbf{PostgreSQL} hautement normalisés avec indexation optimisée et considérations \textbf{Read/Write split} pour absorber les pics de trafic massifs.
\item \textbf{Intégration} de la synchronisation en temps réel de l'inventaire des places avec \textbf{WebSockets (Socket.io)} et transactions atomiques pour éliminer les risques de double réservation.
\item \textbf{Déploiement} via des \textbf{pipelines CI/CD automatisés} avec portes qualité rigoureuses, scénarios de tests de charge et conteneurisation \textbf{Docker}.
\end{highlights}

\textbf{4. Roehrdanz Immobilien --- CMS et Plateforme Immobilière Moderne(\href{https://rimmobilien-test.de}{rimmobilien-test.de})}
\begin{highlights}
\item \textbf{Vue d'ensemble :} Gestion des annonces immobilières, des demandes clients et des présentations de services avec panneau d'administration intégré.
\item \textbf{Pilotage} de la livraison d'une plateforme immobilière moderne avec un \textbf{CMS MERN personnalisé} (React 19, Express 5, MongoDB).
\item \textbf{Développement} d'un panneau d'administration robuste avec \textbf{traitement d'images Base64 asynchrone} et rendu de contenu dynamique optimisé SEO.
\item \textbf{Création} d'un tableau de bord sécurisé réservé aux employés avec \textbf{authentification par session}, hachage Bcryptjs et \textbf{RBAC} pour la modération du contenu.
\item \textbf{Optimisation} des performances frontend avec \textbf{Vite} et \textbf{Tailwind CSS v4}, atteignant des temps de chargement quasi-instantanés pour les galeries photos haute résolution.
\end{highlights}

\textbf{5. B-moovd --- Écosystème Mobile Haute Performance}
\begin{highlights}
\item \textbf{Vue d'ensemble :} Application mobile haute performance pour la réservation fluide, l'orchestration de paiements et l'engagement utilisateur en temps réel.
\item \textbf{Développement} d'une application mobile cross-platform réactive avec \textbf{Flutter}, implémentant une \textbf{gestion d'état par Provider} et des animations fluides à \textbf{60 FPS}.
\item \textbf{Intégration} des appels API RESTful avec \textbf{Dio} et mise en cache persistante locale pour une expérience \textbf{offline-first} fluide.
\item \textbf{Optimisation} de la taille du bundle et de l'allocation des ressources pour des performances élevées sur une large gamme de terminaux mobiles.
\end{highlights}

\textbf{6. B-moovd Ticketing --- Plateforme de Billetterie Événementielle (\href{https://bmoovd-shop.com}{bmoovd-shop.com})}
\begin{highlights}
\item \textbf{Vue d'ensemble :} Plateforme de billetterie événementielle en production pour le B'moovd Sportsbar Wolfsburg, permettant la vente de billets en ligne pour des événements de visionnage sportif.
\item \textbf{Développement} d'une application web de billetterie complète en \textbf{Next.js} pour le venue B'moovd --- filiale du \textbf{Röhrdanz Unternehmensgruppe} --- vendant actuellement des billets pour les événements de visionnage public de la \textbf{Coupe du Monde FIFA 2026} à Wolfsburg.
\item \textbf{Implémentation} des flux de listing d'événements, sélection de billets et paiement avec \textbf{orchestration de paiements complète}.
\item \textbf{Livraison} d'un storefront mobile-optimisé avec une identité visuelle alignée sur la marque B'moovd.
\item \textbf{Déploiement} sous la marque Axynoxia, servant des milliers de fans sur plusieurs événements de la Coupe du Monde.
\end{highlights}

\textbf{7. Confortex --- E-commerce Shopify (\href{https://confortex.shop}{confortex.shop})}
\begin{highlights}
\item \textbf{Développement} d'une boutique \textbf{Shopify} complète pour une marque tunisienne premium de literie, couvrant matelas, oreillers, lits, produits bébé et linge de maison.
\item \textbf{Configuration} d'une architecture produit multi-collections, d'un thème storefront personnalisé et de contenus blog SEO-optimisés pour générer du trafic organique.
\item \textbf{Implémentation} de promotions de livraison gratuite, d'intégration newsletter et d'un système de contact multi-showrooms sur 3 points de vente physiques à Ariana, Tunisie.
\end{highlights}

\textbf{8. Vie Sublime --- E-commerce Shopify (\href{https://viesublime.tn}{viesublime.tn})}
\begin{highlights}
\item \textbf{Création} d'un storefront \textbf{Shopify} premium pour une marque tunisienne haut de gamme spécialisée dans le mobilier artisanal intérieur, bureau, extérieur et scolaire.
\item \textbf{Structuration} d'un catalogue multi-collections (Tables et Chaises, Mobilier Intérieur, Bureau, Scolaire, Jardin) avec un thème personnalisé aligné sur une identité de marque luxe.
\item \textbf{Livraison} d'un blog SEO axé contenu, d'une intégration newsletter et d'une gestion multi-showrooms pour soutenir les canaux de vente en ligne et physiques.
\end{highlights}

\textbf{9. MERN E-commerce --- Infrastructure SaaS Industrielle}
\begin{highlights}
\item \textbf{Vue d'ensemble :} Plateforme retail industrielle pour la gestion des catalogues produits, des commandes clients et des flux de paiement sécurisés.
\item \textbf{Développement} d'un écosystème e-commerce complet avec la stack \textbf{MERN}, axé sur la sécurité industrielle via \textbf{JWT} et \textbf{Argon2}.
\item \textbf{Création} d'un tableau de bord admin avancé avec suivi des commandes en temps réel, gestion des stocks et gestion d'assets cloud intégrée via \textbf{Multer}.
\item \textbf{Implémentation} d'une architecture frontend responsive avec \textbf{Redux Toolkit} pour la persistance d'état complexe dans les workflows d'achat.
\end{highlights}

\textbf{10. OMHY Entertainment --- Hub Vidéo Managé (\href{https://www.omhyfamily.com/}{omhyfamily.com})}
\begin{highlights}
\item \textbf{Vue d'ensemble :} Hub média digital premium pour le streaming vidéo haute fidélité, portfolios de contenu et expériences utilisateur interactives.
\item \textbf{Développement} d'une plateforme de streaming vidéo et de portfolio avec \textbf{React (Chakra UI)} et \textbf{GSAP} pour des animations interactives haute fidélité.
\item \textbf{Orchestration} de déploiements Zero-Downtime via des pipelines \textbf{Jenkins} avec portes qualité automatisées et audits de sécurité.
\item \textbf{Gestion} d'architectures de bases de données hybrides (MongoDB et MySQL) pour équilibrer les métadonnées de contenu flexibles et l'historique strict des transactions utilisateurs.
\end{highlights}

\textbf{11. Jektis Travel --- Plateforme E-Tourisme Mondiale}
\begin{highlights}
\item \textbf{Vue d'ensemble :} Plateforme mondiale de découverte touristique pour la réservation de vols et hôtels avec scraping automatisé de prix en temps réel et paiements sécurisés.
\item \textbf{Architecture} d'une plateforme touristique automatisée en \textbf{PHP}, intégrant \textbf{Stripe} pour l'orchestration des paiements et des modules de scraping pour la découverte en temps réel.
\item \textbf{Développement} de workflows transactionnels automatisés avec \textbf{PHPMailer} et cron jobs personnalisés pour les mises à jour de statut de réservation et les notifications.
\end{highlights}

\textbf{12. New Journey --- Écosystème de Conseil Mondial}
\begin{highlights}
\item \textbf{Vue d'ensemble :} Écosystème de conseil mondial pour connecter clients et experts, gérant les interactions professionnelles et la prestation de services.
\item \textbf{Développement} d'une plateforme de conseil complète avec \textbf{Symfony 7} et \textbf{PHP 8.2}, axée sur une architecture de contrôleurs modulaire et \textbf{Doctrine ORM}.
\item \textbf{Implémentation} de composants UI dynamiques avec \textbf{Symfony UX Turbo} et Stimulus pour une expérience utilisateur haute performance de type SPA.
\item \textbf{Orchestration} de l'environnement de développement avec \textbf{Docker Compose} pour garantir la parité entre développement local et déploiements en production.
\item \textbf{Intégration} d'une authentification multi-utilisateurs sécurisée et d'un contrôle d'accès basé sur les rôles (RBAC) pour gérer les interactions clients et administrateurs.
\end{highlights}

\textbf{13. Hub CI/CD DevOps --- Infrastructure Industrielle (ESPRIT)}
\begin{highlights}
\item \textbf{Construction} d'une plateforme CI/CD complète avec \textbf{Jenkins}, \textbf{SonarQube} et \textbf{Nexus} pour le versionnement des artefacts et la gouvernance de la qualité du code.
\item \textbf{Automatisation} du déploiement de microservices \textbf{Angular} et \textbf{Spring Boot} avec audit de sécurité complet via \textbf{Trivy}.
\end{highlights}

\textbf{14. Lab Cloud Privé OpenStack --- Infrastructure Hybride (ESPRIT)}
\begin{highlights}
\item \textbf{Déploiement} d'un cloud privé \textbf{OpenStack} multi-nœuds, orchestrant réseaux virtuels, sous-réseaux et nœuds \textbf{VPN-over-physical}.
\item \textbf{Gestion} du cycle de vie complet des charges de travail virtualisées, avec mise en œuvre du scaling automatisé et du routage haute disponibilité pour les applications web invitées.
\end{highlights}

% ================= EDUCATION =================
\section{Formation}
\textbf{Diplôme d'Ingénieur en Informatique --- ESPRIT} \hfill 2020 -- 2025 \\
Ariana, Tunisie \\
3 ans Génie Logiciel Général \;|\; 2 ans Spécialité : \textbf{ArcTIC --- Architecture Cloud et DevOps}

% ================= LANGUAGES =================
\section{Langues}
Arabe (Natif) \;|\; Anglais (B2) \;|\; Français (B2)

\end{document}